\chapter{Tổng quan tài liệu và từ khóa}
\label{chap:tong-quan}

\section{Phương pháp tìm tài liệu và bộ từ khóa}
\label{sec:tim-tai-lieu}

Điểm xuất phát là ba bài giảng viên giao cho chương~10 ``Phát hiện lỗi và định vị lỗi'' (Mục~\ref{sec:ba-bai-bao}).
Nhóm lần ngược tài liệu được trích dẫn và tìm xuôi nghiên cứu mới trên arXiv, ACM DL, IEEE Xplore và Google Scholar,
chọn các phương pháp SemLoc so sánh hay kế thừa, bộ dữ liệu lỗi Python, và các bài về cách đánh giá, rủi ro khi dùng LLM,
nhãn sai và bản sửa ``lách'' kiểm thử. Bảng~\ref{tab:tu-khoa} là bộ từ khóa song ngữ dùng để tìm và thống nhất thuật ngữ.

\begin{table}[H]
\centering
\caption{Bộ từ khóa của đề tài}
\label{tab:tu-khoa}
\small
\begin{tabular}{@{}ll@{}}
\toprule
Tiếng Anh & Tiếng Việt \\
\midrule
fault localization & định vị lỗi \\
spectrum-based fault localization (SBFL) & định vị lỗi dựa trên phổ độ phủ \\
semantic fault & lỗi ngữ nghĩa \\
semantic constraint / invariant & ràng buộc ngữ nghĩa / bất biến \\
semantic violation spectrum & phổ vi phạm ràng buộc ngữ nghĩa \\
counterfactual verification & kiểm chứng phản thực tế \\
large language model (LLM) & mô hình ngôn ngữ lớn \\
structured grounding & neo lập luận vào biểu diễn có cấu trúc \\
program instrumentation & chèn mã kiểm tra vào chương trình \\
reproducibility & tính tái hiện được \\
data contamination / memorization & nhiễm dữ liệu huấn luyện / ghi nhớ \\
omission fault & lỗi thiếu mã \\
test-driven development (TDD) & phát triển hướng kiểm thử \\
plausible patch / test overfitting & bản sửa hợp lệ bề ngoài / quá khớp kiểm thử \\
label error (benchmark) & nhãn sai trong bộ dữ liệu \\
static / dynamic analysis & phân tích tĩnh / phân tích động \\
\bottomrule
\end{tabular}
\end{table}

\section{Ba bài báo được giao}
\label{sec:ba-bai-bao}

\textbf{SemLoc}~\cite{yang2026semloc} nhắm vào lỗi ngữ nghĩa, nơi phổ độ phủ ``sụp''. Phương pháp gồm bốn bước: tìm điểm
neo (đầu hàm, định nghĩa -- sử dụng, biên vòng lặp, lệnh trả về); một lượt gọi LLM sinh ràng buộc gắn với điểm neo; chèn
ràng buộc, chạy kiểm thử, chấm ma trận ``ràng buộc $\times$ kiểm thử'' bằng Ochiai; vá thử một dòng để kiểm chứng phản
thực tế. Trên SemFault-250 (250 chương trình Python một hàm, một lỗi) bài báo báo cáo Acc@1
\RqOneSemlocClaudeTZeroEightAuthorAllAccOne\%, nhưng chỉ có kết quả cấp dòng cho bảy lỗi thật. Mã chấm điểm phá hòa theo
thứ tự dòng, loại chương trình chạy lỗi khỏi mẫu số, mỗi cấu hình chạy một lần; đây là đầu vào của RQ1.

\textbf{EffiHolmes}~\cite{yang2026effiholmes} định vị lỗi hiệu năng, loại lỗi không làm kiểm thử thất bại: chạy chương
trình ở hai mức tải, coi hàm có chi phí tăng theo tải là điểm nóng và cho LLM lập luận trên đường gọi tới đó (140 vấn
đề thật). Nhóm học ý tưởng \emph{đối chiếu hai lần thực thi để lọc nhiễu} (nền của CovGate) và khuôn câu hỏi nghiên cứu.

\textbf{Báo cáo sự cố kết hợp LLM}~\cite{medeiros2026crash} gom cụm vết ngăn xếp của hai hệ thống Java công nghiệp, xếp
hạng phương thức đáng nghi rồi đưa vào prompt; trên 38 lỗi thật, cấu hình tốt nhất định vị được 71/100. Bài học: đưa
bằng chứng đã khai thác vào prompt giúp LLM (ý tưởng của biến thể có độ phủ từng dòng), và đầu ra LLM cần được kiểm tra.

\section{Các hướng liên quan, khoảng trống và vị trí của nhóm}
\label{sec:khoang-trong}

Về định vị lỗi, Ochiai~\cite{abreu2009} là mốc SBFL chuẩn, với hai điểm yếu đánh giá đã biết là hòa điểm và lỗi thiếu
mã~\cite{wong2016survey}; gỡ lỗi chia đôi~\cite{zeller2002tse} là nền cho bước kiểm chứng của SemLoc. Trong định vị lỗi
bằng LLM, AutoFL~\cite{kang2024autofl} và FlexFL~\cite{xu2025flexfl} đưa bằng chứng thực thi vào prompt, Wu và cộng
sự~\cite{wu2023llmfl} thấy nhật ký lỗi giúp LLM rõ rệt, còn Agentless~\cite{xia2024agentless} cho thấy quy trình đơn
giản có thể thắng tác tử phức tạp. Về tính hợp lệ, cách chấm~\cite{pearson2017} và nhãn sai~\cite{northcutt2021labels}
đều có thể đảo thứ hạng, bản sửa vượt kiểm thử thường không đúng~\cite{qi2015kali}, và LLM có thể đã ghi nhớ bộ dữ
liệu~\cite{ramos2024memorize}. Bảng~\ref{tab:so-sanh-tai-lieu} đặt ba bài cạnh hướng của nhóm. Khoảng trống còn lại: chưa ai đo phần độ chính xác thật
sự đến từ bằng chứng thực thi; kết quả SemLoc đến từ một lần chạy với cách chấm lạc quan trên nhãn chưa kiểm tra; và
chưa có so sánh với cách làm đơn giản mà lập trình viên vẫn dùng --- một prompt có kiểm thử thất bại. Ba câu hỏi nghiên
cứu ở Chương~\ref{chap:mo-dau} lấp đúng ba khoảng trống này.

\begin{table}[H]
\centering
\caption{So sánh ba bài báo được giao với hướng của nhóm}
\label{tab:so-sanh-tai-lieu}
\small
\begin{tabularx}{\linewidth}{@{}>{\raggedright\arraybackslash}p{0.17\linewidth}*{4}{>{\raggedright\arraybackslash}X}@{}}
\toprule
Tiêu chí & SemLoc~\cite{yang2026semloc} & EffiHolmes~\cite{yang2026effiholmes} & Báo cáo sự cố + LLM~\cite{medeiros2026crash} & Nhóm M9 \\
\midrule
Loại lỗi & ngữ nghĩa & hiệu năng & sự cố (Java) & ngữ nghĩa \\
Tín hiệu & ràng buộc do LLM sinh + chạy kiểm thử & phân tích hiệu năng ở hai mức tải & vết ngăn xếp + xếp hạng & kiểm thử thất bại đưa thẳng cho LLM; tách đóng góp của từng tín hiệu \\
Lỗi thật & 7 lỗi có kết quả cấp dòng & 140 vấn đề & 38 lỗi công nghiệp & 31 lỗi BugsInPy, 3 lần chạy \\
Kiểm tra bộ dữ liệu & không & --- & --- & kiểm chứng bằng thực thi và rà soát toàn bộ nhãn \\
Lặp và kiểm định & 1 lần, không kiểm định & --- & --- & 3 lần, McNemar, bootstrap theo cụm \\
\bottomrule
\end{tabularx}
\end{table}
